\documentclass[12pt]{article}
\usepackage{amsmath}
\usepackage{amsfonts}
\usepackage{float}
\usepackage{fancyhdr}
\usepackage{graphicx}
\usepackage{setspace}
\usepackage[colorlinks=true,linkcolor=blue, citecolor=red]{hyperref}
\usepackage{url}
\usepackage[top=.75in, left=.75in, right=.75in, bottom=1in]{geometry}

% For algorithm
\usepackage{algorithm}
\usepackage{algpseudocode}

% ============ CODE ============
\usepackage{listings}
\usepackage{xcolor}
\definecolor{codegreen}{rgb}{0,0.6,0}
\definecolor{codegray}{rgb}{0.5,0.5,0.5}
\definecolor{codepurple}{rgb}{0.58,0,0.82}
\definecolor{backcolour}{rgb}{0.95,0.95,0.92}

% Styling for the code.
\lstdefinestyle{mystyle}{
    backgroundcolor=\color{backcolour},   
    commentstyle=\color{codegreen},
    keywordstyle=\color{magenta},
    numberstyle=\tiny\color{codegray},
    stringstyle=\color{codepurple},
    basicstyle=\ttfamily\footnotesize,
    breakatwhitespace=false,         
    breaklines=true,                 
    captionpos=b,                    
    keepspaces=true,                 
    numbers=left,                    
    numbersep=5pt,                  
    showspaces=false,                
    showstringspaces=false,
    showtabs=false,                  
    tabsize=2
}
\lstset{style=mystyle}

% Table packages and Unicode / Font support
\usepackage{array}
\usepackage{longtable}
\usepackage{enumitem}
\usepackage{booktabs}
\usepackage[utf8]{inputenc}
\DeclareUnicodeCharacter{2264}{\ensuremath{\le}}
\DeclareUnicodeCharacter{2194}{\ensuremath{\leftrightarrow}}
\usepackage[T1]{fontenc}
\usepackage{lmodern}
\graphicspath{{../../../../templates/}{../../prototype/}{figures/}{./}}
\setlength{\emergencystretch}{3em}
\setlist[enumerate]{label=\arabic*.,leftmargin=*}
\setlist[itemize]{leftmargin=*}

% ============ HEADER AND FOOTER ============
\setlength{\headheight}{29.43912pt}

\pagestyle{fancy}
\fancyfoot{}
\fancyfoot[R]{Page \thepage}

\setlength{\parindent}{0pt}
\renewcommand{\baselinestretch}{1.5}

\newcommand{\coursename}{Introduction of Software Engineering}
\newcommand{\reportname}{Studitopia}
\newcommand{\reporttitle}{\href{https://github.com/RoundFloor/Studitopia.git}{Studitopia}}

\newcommand{\studentname}{%
Nguyen Thanh An (24125024)\\
Nguyen Chi Bao (24125025)\\
Tran Nguyen Khanh Binh (24125026)\\
Tran Canh Anh Tuan (24125048)\\
Nguyen Le Quoc Huy (24125100)
}
\newcommand{\teachername}{Ho Tuan Thanh, Nguyen Thi Minh Tuyen}

% Header
\lhead{\reporttitle}
\rhead{
University of Science - VNUHCM\\
\coursename
}

% Footer
\newcommand{\leftfooter}{}
\lfoot{\leftfooter}

% ============ DOCUMENT ============
\begin{document}

\pagenumbering{roman}
\begin{titlepage}
\newcommand{\HRule}{\rule{\linewidth}{0.5mm}}
\centering

\textsc{\LARGE VNU-HCM}\\[0.5cm]
\textsc{\Large University of Science}\\[0.5cm]
\textsc{\large Faculty of Information Technology}\\[0.5cm]
\textsc{Computer Science}\\[0.5cm]

\HRule \\[0.4cm]
{ 
\huge{\bfseries{\reporttitle}}\\[0.5cm]
\large{\bfseries{Post Management --- Use Case Diagram}}
}\\[0.4cm]
\HRule \\[0.5cm]

\textbf{\large Subject: \coursename}\\[0.5cm]

\begin{minipage}[t]{0.55\textwidth}
\begin{flushleft} \large
\emph{Students:}\\
\studentname
\end{flushleft}
\end{minipage}
~
\begin{minipage}[t]{0.4\textwidth}
\begin{flushright} \large
\emph{Supervising Teachers:} \\
\teachername
\end{flushright}
\end{minipage}\\[1cm]

{\large \today}\\[1cm]

\includegraphics[scale=.20]{img/hcmus-logo.png}\\[1cm] 

\vfill
\end{titlepage}	

\tableofcontents
\pagebreak

\pagenumbering{arabic}
\setcounter{page}{1}

\hypertarget{sec:post-management-use-case-diagram}{}
\addcontentsline{toc}{section}{Post Management --- Use Case Diagram}
\section*{Post Management --- Use Case Diagram}
\par\medskip

\textbf{Version:} \texttt{1.0}
\par\medskip

\textbf{Performed by:} [Nguyen Thanh An] | \textbf{Edited by:} [Nguyen Thanh An]
\par\medskip

Merged from the focused use case diagrams in this functional group.
\par\medskip

\begin{center}
\includegraphics[width=0.85\textwidth,height=0.75\textheight,keepaspectratio]{"figures/diagram_1.png"}
\end{center}
\par\medskip

\textbf{Diagram notation.} Person shape = human actor; blue = base use case; purple = included; green = extending; solid line = association; dashed arrow = include or extend.
\par\medskip

\pagebreak
\hypertarget{sec:uc-04-create-discussion-post}{}
\addcontentsline{toc}{section}{UC-04 --- Create Discussion Post}
\section*{UC-04 --- Create Discussion Post}
\par\medskip

\textbf{Version:} \texttt{2.0}
\par\medskip

\textbf{Performed by:} [Tran Canh Anh Tuan] | \textbf{Edited by:} [Nguyen Thanh An]
\par\medskip

\begingroup
\small
\renewcommand{\baselinestretch}{1.0}\selectfont
\renewcommand{\arraystretch}{1.15}
\setlength{\tabcolsep}{6pt}
\begin{longtable}{|>{\raggedright\arraybackslash}p{0.09\textwidth}|>{\raggedright\arraybackslash}p{0.84\textwidth}|}
\hline
\textbf{Field} & \textbf{Value} \\ \hline
\endhead
Use Case ID & UC-04 \\ \hline
Name & Create Discussion Post \\ \hline
Primary Actor & University Student \\ \hline
Supporting Actors & None \\ \hline
Feature Group & Forum (Primary) \\ \hline
Priority & High \\ \hline
Related User Stories & US-14 \\ \hline
Related Requirements & US-14, US-17, US-19, US-20; FR-14.1 to FR-14.6, FR-17.1 to FR-17.5, FR-19.1 to FR-19.5, FR-20.1 to FR-20.5; NFR-PERF-4, NFR-SEC-8, NFR-SEC-9 \\ \hline
Source Documents & \texttt{proposals/use\_cases\_model/UC-04-create-discussion-post.md}; \texttt{proposals/use\_cases\_description.md} (UC-04); \texttt{proposals/use\_cases\_model.md}; \texttt{proposals/use\_case\_diagrams.md}; \texttt{PA3-Group05/use\_case.md} (section 4); \texttt{proposals/overview.md}; \texttt{proposals/functional\_requirement.md}; \texttt{proposals/non\_functional\_requirement.md}; \texttt{proposals/userstory.md}; \texttt{proposals/feature\_flow.md}; \texttt{proposals/screen\_flow.md}; \texttt{proposals/API.md}; \texttt{proposals/response\_code.md}; \texttt{proposals/create\_post.md}; \texttt{proposals/list\_post.md} \\ \hline
\end{longtable}
\endgroup
\par\medskip

\hypertarget{sec:part-a-use-case-model}{}
\addcontentsline{toc}{subsection}{Part A --- Use Case Model}
\subsection*{Part A --- Use Case Model}
\par\medskip

\hypertarget{sec:a1-focused-use-case-diagram}{}
\addcontentsline{toc}{subsubsection}{A1. Focused Use Case Diagram}
\subsubsection*{A1. Focused Use Case Diagram}
\par\medskip

\begin{center}
\includegraphics[width=0.85\textwidth,height=0.75\textheight,keepaspectratio]{"figures/UC-04-create-discussion-post_1.png"}
\end{center}
\par\medskip

\textbf{Diagram notation.} Person shape = human actor; blue = selected base use case; purple = included; green = extending; solid line = association; dashed arrow = include or extend.
\par\medskip

\hypertarget{sec:a2-actor-register}{}
\addcontentsline{toc}{subsubsection}{A2. Actor Register}
\subsubsection*{A2. Actor Register}
\par\medskip

\begingroup
\small
\renewcommand{\baselinestretch}{1.0}\selectfont
\renewcommand{\arraystretch}{1.15}
\setlength{\tabcolsep}{6pt}
\begin{longtable}{|>{\raggedright\arraybackslash}p{0.13\textwidth}|>{\raggedright\arraybackslash}p{0.18\textwidth}|>{\raggedright\arraybackslash}p{0.61\textwidth}|}
\hline
\textbf{Actor} & \textbf{Type} & \textbf{Description} \\ \hline
\endhead
University Student & Primary human actor & A registered undergraduate student who publishes a discussion post in the selected university, major, and course forum. \\ \hline
New University Student & Specialized human actor & Inherits the general University Student association to UC-04 from the shared actor model; not drawn separately in the focused diagram. \\ \hline
Registered University Student & Specialized human actor & Inherits the general University Student association to UC-04 from the shared actor model. \\ \hline
Post Author & Secondary human role / specialized actor & The acting University Student role during this use case; inherits the UC-04 association and becomes the author of the stored post. \\ \hline
Room Owner or Manager & Specialized human actor & Inherits the general University Student association to UC-04 from the shared actor model. \\ \hline
\end{longtable}
\endgroup
\par\medskip

\hypertarget{sec:a3-relationships}{}
\addcontentsline{toc}{subsubsection}{A3. Relationships}
\subsubsection*{A3. Relationships}
\par\medskip

\hypertarget{sec:a3-1-associations}{}
\paragraph*{A3.1 Associations}
\par\medskip

\begingroup
\small
\renewcommand{\baselinestretch}{1.0}\selectfont
\renewcommand{\arraystretch}{1.15}
\setlength{\tabcolsep}{6pt}
\begin{longtable}{|>{\raggedright\arraybackslash}p{0.15\textwidth}|>{\raggedright\arraybackslash}p{0.23\textwidth}|>{\raggedright\arraybackslash}p{0.10\textwidth}|>{\raggedright\arraybackslash}p{0.43\textwidth}|}
\hline
\textbf{Actor} & \textbf{Use Case} & \textbf{Type} & \textbf{Description} \\ \hline
\endhead
University Student & UC-04 Create Discussion Post & Primary & Publishes a discussion post in the university forum. \\ \hline
\end{longtable}
\endgroup
\par\medskip

\hypertarget{sec:a3-2-generalization-and-specialization}{}
\paragraph*{A3.2 Generalization and Specialization}
\par\medskip

\begingroup
\small
\renewcommand{\baselinestretch}{1.0}\selectfont
\renewcommand{\arraystretch}{1.15}
\setlength{\tabcolsep}{6pt}
\begin{longtable}{|>{\raggedright\arraybackslash}p{0.42\textwidth}|>{\raggedright\arraybackslash}p{0.10\textwidth}|>{\raggedright\arraybackslash}p{0.10\textwidth}|>{\raggedright\arraybackslash}p{0.29\textwidth}|}
\hline
\textbf{Child} & \textbf{Parent} & \textbf{Type} & \textbf{Inherited Behavior or Association} \\ \hline
\endhead
New University Student; Registered University Student; Post Author; Room Owner or Manager & University Student & Actor generalization & Each specialized student actor inherits the UC-04 association. \\ \hline
\end{longtable}
\endgroup
\par\medskip

\hypertarget{sec:a3-3-include}{}
\paragraph*{A3.3 Include}
\par\medskip

\begingroup
\small
\renewcommand{\baselinestretch}{1.0}\selectfont
\renewcommand{\arraystretch}{1.15}
\setlength{\tabcolsep}{6pt}
\begin{longtable}{|>{\raggedright\arraybackslash}p{0.17\textwidth}|>{\raggedright\arraybackslash}p{0.16\textwidth}|>{\raggedright\arraybackslash}p{0.59\textwidth}|}
\hline
\textbf{Base Use Case} & \textbf{Included Use Case} & \textbf{Meaning} \\ \hline
\endhead
UC-04 Create Discussion Post & I-05 Validate Post Content & Validate title, content, tags, post type, image support, and forum context before publication. \\ \hline
\end{longtable}
\endgroup
\par\medskip

\hypertarget{sec:a3-4-extend}{}
\paragraph*{A3.4 Extend}
\par\medskip

\begingroup
\small
\renewcommand{\baselinestretch}{1.0}\selectfont
\renewcommand{\arraystretch}{1.15}
\setlength{\tabcolsep}{6pt}
\begin{longtable}{|>{\raggedright\arraybackslash}p{0.16\textwidth}|>{\raggedright\arraybackslash}p{0.18\textwidth}|>{\raggedright\arraybackslash}p{0.25\textwidth}|>{\raggedright\arraybackslash}p{0.32\textwidth}|}
\hline
\textbf{Extending Use Case} & \textbf{Base Use Case} & \textbf{Condition} & \textbf{Added Behavior} \\ \hline
\endhead
E-05 Attach Image or LaTeX & UC-04 Create Discussion Post & The author adds rich content to a post. & Add image and LaTeX content to the discussion post. \\ \hline
\end{longtable}
\endgroup
\par\medskip

\hypertarget{sec:a4-model-notes}{}
\addcontentsline{toc}{subsubsection}{A4. Model Notes}
\subsubsection*{A4. Model Notes}
\par\medskip

\begin{enumerate}
\item Every discussion post must include I-05 because invalid or incomplete post content cannot be published.
\item Image and LaTeX content are optional extensions of the base post-creation behavior through E-05.
\item Post Author is the acting University Student role during this use case; the general UC-04 association is inherited by all specialized student actors.
\item The focused diagram matches the UC-04 diagram in \texttt{use\_case\_diagrams.md} and contains no supporting system actor or specialized use case.
\item Actor generalization is documented in the focused model relationship summary and the shared actor model even though the focused diagram shows only the general University Student association.
\end{enumerate}
\par\medskip

\hypertarget{sec:part-b-use-case-specification}{}
\addcontentsline{toc}{subsection}{Part B --- Use Case Specification}
\subsection*{Part B --- Use Case Specification}
\par\medskip

\hypertarget{sec:b1-brief-description}{}
\addcontentsline{toc}{subsubsection}{B1. Brief Description}
\subsubsection*{B1. Brief Description}
\par\medskip

A student publishes an academic question or general discussion in the selected university, major, and course forum. The post supports a title, text content, optional LaTeX, optional images, and tags. The trigger is the student selecting \textbf{Create Post} and choosing the discussion-post path; the successful outcome is a validated discussion post published in the selected forum context.
\par\medskip

\hypertarget{sec:b2-preconditions}{}
\addcontentsline{toc}{subsubsection}{B2. Preconditions}
\subsubsection*{B2. Preconditions}
\par\medskip

\begin{enumerate}
\item The student is authenticated with an active session.
\item The student has selected a valid university, major, and course forum context.
\item The student is permitted to create posts in that forum.
\end{enumerate}
\par\medskip

\hypertarget{sec:b3-postconditions}{}
\addcontentsline{toc}{subsubsection}{B3. Postconditions}
\subsubsection*{B3. Postconditions}
\par\medskip

\begin{enumerate}
\item \textbf{Success:} A validated discussion post is stored with its author, content, post type, tags, and forum context.
\item \textbf{Success:} The post appears in the corresponding forum post list and can be opened in post detail.
\item \textbf{Failure:} No post is published, and the student receives a clear validation or system error with the opportunity to correct the content.
\end{enumerate}
\par\medskip

\hypertarget{sec:b4-main-success-scenario-basic-flow}{}
\addcontentsline{toc}{subsubsection}{B4. Main Success Scenario (Basic Flow)}
\subsubsection*{B4. Main Success Scenario (Basic Flow)}
\par\medskip

\begin{enumerate}
\item The student selects \textbf{Create Post} from the selected forum.
\item The system displays the post-type choices and the discussion post editor.
\item The system shows the active university, major, and course context for the new post.
\item The student enters a title, text content, optional LaTeX, optional images, and tags.
\item The student marks the post as a question or a general discussion.
\item The frontend validates the required content and supported image types or sizes before submission.
\item The student submits the post.
\item The backend verifies authentication, forum access, post permissions, content, and rate limits.
\item The backend stores the post with its author, post type, tags, and university, major, and course associations.
\item The system publishes the post and returns the new post result.
\item The system displays the post in the corresponding forum post list and allows the student to open its detail page.
\end{enumerate}
\par\medskip

\hypertarget{sec:b5-alternative-flows}{}
\addcontentsline{toc}{subsubsection}{B5. Alternative Flows}
\subsubsection*{B5. Alternative Flows}
\par\medskip

Unless otherwise stated, an alternative flow resumes at the next step of the basic flow.
\par\medskip

\hypertarget{sec:a1-general-discussion}{}
\paragraph*{A1. General discussion}
\par\medskip

At basic flow step 5, the student selects \textbf{General Discussion} instead of \textbf{Question}. The publication flow is unchanged, but the post is not treated as a question awaiting a true response.
\par\medskip

\hypertarget{sec:a2-question-post}{}
\paragraph*{A2. Question post}
\par\medskip

At basic flow step 5, the student selects \textbf{Question}. The post becomes eligible for answers and true-response marking after publication.
\par\medskip

\hypertarget{sec:a3-add-or-remove-optional-content}{}
\paragraph*{A3. Add or remove optional content}
\par\medskip

At basic flow step 4, the student adds or removes LaTeX, images, or tags. The editor updates the preview and validation state before submission.
\par\medskip

\hypertarget{sec:a4-cancel-creation}{}
\paragraph*{A4. Cancel creation}
\par\medskip

The student leaves the editor before submitting. No published post is created; any retained draft behavior is subject to the platform's buffering policy.
\par\medskip

\hypertarget{sec:b6-exception-flows}{}
\addcontentsline{toc}{subsubsection}{B6. Exception Flows}
\subsubsection*{B6. Exception Flows}
\par\medskip

\hypertarget{sec:e1-missing-or-invalid-content}{}
\paragraph*{E1. Missing or invalid content}
\par\medskip

The system highlights the invalid fields, keeps the editor open, and does not submit until corrected.
\par\medskip

\hypertarget{sec:e2-unsupported-or-oversized-image}{}
\paragraph*{E2. Unsupported or oversized image}
\par\medskip

The system rejects the image, identifies the constraint, and allows the student to choose another file.
\par\medskip

\hypertarget{sec:e3-forum-access-lost}{}
\paragraph*{E3. Forum access lost}
\par\medskip

If the forum context is no longer valid or the student lacks posting permission, the system rejects publication and explains the reason.
\par\medskip

\hypertarget{sec:e4-rate-limit-reached}{}
\paragraph*{E4. Rate limit reached}
\par\medskip

The system rejects the submission and informs the student to retry later.
\par\medskip

\hypertarget{sec:e5-server-or-database-failure}{}
\paragraph*{E5. Server or database failure}
\par\medskip

No post is published. The system displays a retry option and preserves the student's entered content where supported.
\par\medskip

\hypertarget{sec:b7-business-rules}{}
\addcontentsline{toc}{subsubsection}{B7. Business Rules}
\subsubsection*{B7. Business Rules}
\par\medskip

\begin{enumerate}
\item A discussion post must be associated with the selected university, major, and course.
\item A discussion post supports a title, text content, LaTeX, images, and tags.
\item The author must choose whether the post is a question or a general discussion.
\item Content must pass frontend and backend validation before publication.
\item A published post must appear in the corresponding forum post list.
\item Only authenticated and authorized students may create posts in a forum.
\item Post publication is subject to rate limiting to reduce spam and abuse.
\end{enumerate}
\par\medskip

\hypertarget{sec:b8-special-requirements}{}
\addcontentsline{toc}{subsubsection}{B8. Special Requirements}
\subsubsection*{B8. Special Requirements}
\par\medskip

\begingroup
\small
\renewcommand{\baselinestretch}{1.0}\selectfont
\renewcommand{\arraystretch}{1.15}
\setlength{\tabcolsep}{6pt}
\begin{longtable}{|>{\raggedright\arraybackslash}p{0.11\textwidth}|>{\raggedright\arraybackslash}p{0.82\textwidth}|}
\hline
\textbf{ID} & \textbf{Requirement} \\ \hline
\endhead
NFR-PERF-4 & A new post must be reflected in the relevant forum view within 2 seconds. \\ \hline
NFR-SEC-8 & The system must protect against SQL injection, XSS, CSRF, and other common web vulnerabilities. \\ \hline
NFR-SEC-9 & Post creation must be rate-limited to reduce spam and abuse. \\ \hline
Design constraint & The editor must validate required content, supported image types, and image sizes before submission. LaTeX and images are optional. \\ \hline
\end{longtable}
\endgroup
\par\medskip

\hypertarget{sec:b9-extension-points}{}
\addcontentsline{toc}{subsubsection}{B9. Extension Points}
\subsubsection*{B9. Extension Points}
\par\medskip

\begingroup
\small
\renewcommand{\baselinestretch}{1.0}\selectfont
\renewcommand{\arraystretch}{1.15}
\setlength{\tabcolsep}{6pt}
\begin{longtable}{|>{\raggedright\arraybackslash}p{0.42\textwidth}|>{\raggedright\arraybackslash}p{0.51\textwidth}|}
\hline
\textbf{Extension Point} & \textbf{Location in Flow of Events} \\ \hline
\endhead
Select post type & Basic flow step 5; alternatives A1-A2 \\ \hline
Add optional rich content & Basic flow step 4; alternative A3 \\ \hline
Image validation failure & Basic flow step 6; exception E2 \\ \hline
Publication authorization failure & Basic flow step 8; exception E3 \\ \hline
Rate-limit recovery & Basic flow step 8; exception E4 \\ \hline
Cancel creation & Before basic flow step 7; alternative A4 \\ \hline
\end{longtable}
\endgroup
\par\medskip

\hypertarget{sec:b10-interface-and-data-notes}{}
\addcontentsline{toc}{subsubsection}{B10. Interface and Data Notes}
\subsubsection*{B10. Interface and Data Notes}
\par\medskip

\begin{itemize}
\item \textbf{Screen/flow path:} \texttt{screen\_flow.md} gives \texttt{Forum -$>$ University -$>$ Major -$>$ Course -$>$ Post} and \texttt{Forum -$>$ Create Post -$>$ Post}. \texttt{feature\_flow.md} gives \texttt{Create Post -$>$ Discussion Post / Material Post / Functional Post}, with discussion content branches \texttt{Text / LaTeX}, \texttt{Image}, \texttt{Interactive Box}, and \texttt{Title}, and all post types converging on \textbf{Submit}.
\item \textbf{Editor fields and post-type choices:} The shared Create Post page offers \textbf{Discussion Post}, \textbf{Material Post}, and \textbf{Functional Post}; UC-04 covers the Discussion Post editor. Discussion fields are title, text content, optional LaTeX, optional images, and tags. The discussion category is \textbf{Question} or \textbf{General Discussion}. \texttt{create\_post.md} recommends the component names \texttt{CreatePostPage}, \texttt{PostTypeSelector}, \texttt{DiscussionEditorForm}, \texttt{PostCategorySelector}, \texttt{DiscussionContentEditor}, \texttt{LatexInsertButton}, \texttt{ImageUploadButton}, \texttt{PostPreviewPanel}, \texttt{CreatePostValidationSummary}, and \texttt{SubmitPostButton}.
\item \textbf{Context data:} The post is associated with \texttt{university\_id}, optional \texttt{major\_id}, and optional \texttt{course\_id}, and the editor shows the active university, major, and course context (\texttt{FR-14.4}, \texttt{AC-14.3}, UC-04 step 3).
\item \textbf{Validation, permission, and rate-limit rules:} Required content and supported image types/sizes are validated before submission (\texttt{FR-14.5}, design constraint). Backend checks authentication, forum access, posting permission, content, and rate limits (\texttt{NFR-SEC-6}, \texttt{NFR-SEC-8}, \texttt{NFR-SEC-9}, UC-04 step 8 and business rule 7). Concrete image formats and size limits are not specified in the sources.
\item \textbf{Exact UI messages available in the sources:} \texttt{create\_post.md} recommends \texttt{ForumAccessDeniedPanel} --- "You do not have permission to post in this forum."; \texttt{RateLimitError} --- "Too many posts. Please retry later."; \texttt{ServerFailureAlert} --- "Post could not be published."; \texttt{CreatePostSuccessToast} --- "Post published."; \texttt{UnsavedChangesDialog} --- "Discard unsaved changes?"; \texttt{DraftRestoreNotice} --- "A saved draft was restored."; \texttt{PostTitleInput} helper "Required"; \texttt{CreatePostButton} label "Create Post"; and \texttt{SubmitPostButton} label "Submit Post" or "Publish Post". These are recommended labels/copy, not mandated requirement text.
\item \textbf{API endpoint:} Backend route is \texttt{POST /forum/posts}. Body fields are \texttt{university\_id: string}, \texttt{major\_id?: string}, \texttt{course\_id?: string}, \texttt{post\_type: 'discussion' | 'material' | 'functional'}, \texttt{title: string}, \texttt{content: string}, \texttt{tags?: string[]}, and \texttt{files?: File[]}. Success is \texttt{201} with \texttt{data: \{ post\_id: string, created\_at: number \}}. Failures are \texttt{400} with \texttt{(INVALID\_POST\_TYPE | INVALID\_POST\_CONTENT)[]}, \texttt{401}, \texttt{403} with \texttt{(FORBIDDEN)[]}, and \texttt{500} with \texttt{(DATABASE\_ERROR | UNKNOWN\_ERROR)[]}.
\item \textbf{Database contract:} \texttt{studitopia.create\_forum\_post(pauthor\_id: uuid, puniversity\_id: uuid, pmajor\_id: uuid, pcourse\_id: uuid, ppost\_type: string, ptitle: string, pcontent: string, ptags: string[], pfiles: jsonb)} returns \texttt{\{post\_id: uuid, created\_at: timestamp\}}. Declared return errors are \texttt{UNKNOWN\_ERROR: 45000}, \texttt{INVALID\_POST\_TYPE: 48003}, \texttt{INVALID\_POST\_CONTENT: 48004}, and \texttt{FORBIDDEN: 48010}.
\item \textbf{Response-code register:} Forum codes include \texttt{48003} \texttt{INVALID\_POST\_TYPE}, \texttt{48004} \texttt{INVALID\_POST\_CONTENT}, and \texttt{48010} \texttt{FORBIDDEN}; backend codes also carry \texttt{45000} \texttt{UNKNOWN DATABASE ERROR}. No dedicated rate-limit status or error code is defined for post creation in \texttt{API.md} or \texttt{response\_code.md}.
\item \textbf{Post-publication surface:} The new post must appear in the corresponding forum post list; \texttt{list\_post.md} defines the post summary as title, author, post type, timestamp, vote score, comment count, solved status, and tags, and opens \texttt{PostDetailPage} with full content, metadata, discussion, attachments, and tools where present.
\item \textbf{Functional-post extra detail:} Outside the UC-04 discussion path but on the same Create Post page, \texttt{US-16} / \texttt{FR-16.1} to \texttt{FR-16.6} / \texttt{AC-16.1} to \texttt{AC-16.5} define an interactive academic tool: an immediate quiz with multiple-choice or single-choice questions, a prerequisite tree created manually or generated from a source file, or dynamic computation. \texttt{create\_post.md} recommends drag-and-drop and code authoring modes, \texttt{Auto apply}, \texttt{Simulate by click}, and mapping invalid tool content to generic \texttt{INVALID\_POST\_CONTENT} because \texttt{API.md} has no tree-specific error contract.
\end{itemize}
\par\medskip

\hypertarget{sec:b11-source-alignment-notes}{}
\addcontentsline{toc}{subsubsection}{B11. Source Alignment Notes}
\subsubsection*{B11. Source Alignment Notes}
\par\medskip

\begin{itemize}
\item \texttt{use\_cases\_description.md} and the PA3 specification describe UC-04 as discussion-post creation only, while \texttt{overview.md}, \texttt{feature\_flow.md}, and \texttt{create\_post.md} place three post types on a shared Create Post page; material- and functional-post details are recorded only as interface context, not folded into the UC-04 basic flow.
\item The focused diagram shows only the general University Student actor, whereas \texttt{use\_cases\_model.md} and the focused model relationship table define actor generalization for specialized student actors; the inheritance is documented in A2/A3 rather than invented in the diagram.
\item \texttt{API.md} uses one \texttt{post\_type} enum for \texttt{discussion | material | functional}, but UC-04 and \texttt{FR-14.1} to \texttt{FR-14.6} constrain this use case to the discussion path.
\item \texttt{API.md} declares \texttt{files?: File[]} at the backend boundary and \texttt{pfiles: jsonb} at the database boundary; the sources do not define a separate image-only schema for UC-04 despite the editor-level image requirement.
\item Rate limiting is required by UC-04 business rule 7 and \texttt{NFR-SEC-9}, but \texttt{API.md} and \texttt{response\_code.md} define no \texttt{429} response or rate-limit error code; exception E4 and the recommended \texttt{RateLimitError} message are the only documented handling.
\item Concrete supported image formats, maximum image sizes, and upload limits are absent from the sources; \texttt{create\_post.md} deliberately avoids naming formats or limits.
\item UI routes, component names, and quoted messages in \texttt{create\_post.md} and \texttt{list\_post.md} are explicitly recommendations derived from the proposals; the requirement documents do not mandate those exact strings.
\item \texttt{FR-17}, \texttt{FR-19}, and \texttt{FR-20} are listed as related because a published Question post becomes eligible for detail, voting, answers, and true-response marking; those downstream actions belong to UC-03/UC-05 and US-19/US-20, not to the publication steps of UC-04.
\item \texttt{FR-04.1} states a 15-minute inactivity timeout while the PA3 UC-02 special requirements state 30 minutes; this conflict concerns session management and is not imported into UC-04 special requirements.
\end{itemize}
\par\medskip

\par\medskip
\hypertarget{sec:ui-prototype-specimens}{}
\addcontentsline{toc}{subsection}{UI Prototype Specimens}
\subsection*{UI Prototype Specimens}
\par\medskip
\begin{figure}[H]
\centering
\includegraphics[width=0.75\textwidth,height=0.42\textheight,keepaspectratio]{"../../prototype/create_post_desktop.png"}
\par\vspace{2pt}
{\footnotesize \textbf{UI Prototype:} Create post - Desktop}
\end{figure}
\par\smallskip
\begin{figure}[H]
\centering
\includegraphics[width=0.75\textwidth,height=0.42\textheight,keepaspectratio]{"../../prototype/create_post_discussion.png"}
\par\vspace{2pt}
{\footnotesize \textbf{UI Prototype:} Create post - Discussion}
\end{figure}
\par\smallskip
\begin{figure}[H]
\centering
\includegraphics[width=0.75\textwidth,height=0.42\textheight,keepaspectratio]{"../../prototype/create_post_validation_and_publication_states.png"}
\par\vspace{2pt}
{\footnotesize \textbf{UI Prototype:} Create post - Validation and publication states}
\end{figure}
\par\smallskip

\pagebreak
\hypertarget{sec:uc-15-create-material-post}{}
\addcontentsline{toc}{section}{UC-15 --- Create Material Post}
\section*{UC-15 --- Create Material Post}
\par\medskip

\textbf{Version:} \texttt{1.0}
\par\medskip

\textbf{Performed by:} [Nguyen Thanh An] | \textbf{Edited by:} [Nguyen Thanh An]
\par\medskip

\begingroup
\small
\renewcommand{\baselinestretch}{1.0}\selectfont
\renewcommand{\arraystretch}{1.15}
\setlength{\tabcolsep}{6pt}
\begin{longtable}{|>{\raggedright\arraybackslash}p{0.09\textwidth}|>{\raggedright\arraybackslash}p{0.84\textwidth}|}
\hline
\textbf{Field} & \textbf{Value} \\ \hline
\endhead
Use Case ID & UC-15 \\ \hline
Name & Create Material Post \\ \hline
Primary Actor & University Student \\ \hline
Supporting Actors & None \\ \hline
Feature Group & Post Management \\ \hline
Priority & Not explicitly stated \\ \hline
Related User Stories & US-15 \\ \hline
Related Requirements & FR-15.1 to FR-15.6; NFR-PERF-4, NFR-SEC-1, NFR-SEC-5, NFR-SEC-6, NFR-SEC-9, NFR-PRIV-5, NFR-PRIV-7, NFR-SCAL-4 \\ \hline
Source Documents & \texttt{proposals/userstory.md} (section 5.2); \texttt{proposals/functional\_requirement.md} (US-15); \texttt{proposals/create\_post.md}; \texttt{proposals/list\_post.md}; \texttt{proposals/post\_management.md}; \texttt{proposals/API.md}; \texttt{proposals/response\_code.md}; \texttt{proposals/non\_functional\_requirement.md}; \texttt{proposals/feature\_flow.md}; \texttt{proposals/screen\_flow.md} \\ \hline
\end{longtable}
\endgroup
\par\medskip

\hypertarget{sec:part-a-use-case-model}{}
\addcontentsline{toc}{subsection}{Part A --- Use Case Model}
\subsection*{Part A --- Use Case Model}
\par\medskip

\hypertarget{sec:a1-focused-use-case-diagram}{}
\addcontentsline{toc}{subsubsection}{A1. Focused Use Case Diagram}
\subsubsection*{A1. Focused Use Case Diagram}
\par\medskip

\begin{center}
\includegraphics[width=0.85\textwidth,height=0.75\textheight,keepaspectratio]{"figures/UC-15-create-material-post_1.png"}
\end{center}
\par\medskip

\textbf{Diagram notation.} Person shape = human actor; cloud = supporting system; blue = selected use case; solid line = association; labelled solid arrow = generalization; dashed arrow = include or extend.
\par\medskip

\hypertarget{sec:a2-actor-register}{}
\addcontentsline{toc}{subsubsection}{A2. Actor Register}
\subsubsection*{A2. Actor Register}
\par\medskip

\begingroup
\small
\renewcommand{\baselinestretch}{1.0}\selectfont
\renewcommand{\arraystretch}{1.15}
\setlength{\tabcolsep}{6pt}
\begin{longtable}{|>{\raggedright\arraybackslash}p{0.10\textwidth}|>{\raggedright\arraybackslash}p{0.10\textwidth}|>{\raggedright\arraybackslash}p{0.72\textwidth}|}
\hline
\textbf{Actor} & \textbf{Type} & \textbf{Description} \\ \hline
\endhead
University Student & Primary human actor & A university student who creates a material post in the selected forum so that other students can access shared learning resources. \\ \hline
\end{longtable}
\endgroup
\par\medskip

\hypertarget{sec:a3-relationships}{}
\addcontentsline{toc}{subsubsection}{A3. Relationships}
\subsubsection*{A3. Relationships}
\par\medskip

\hypertarget{sec:a3-1-associations}{}
\paragraph*{A3.1 Associations}
\par\medskip

\begingroup
\small
\renewcommand{\baselinestretch}{1.0}\selectfont
\renewcommand{\arraystretch}{1.15}
\setlength{\tabcolsep}{6pt}
\begin{longtable}{|>{\raggedright\arraybackslash}p{0.11\textwidth}|>{\raggedright\arraybackslash}p{0.16\textwidth}|>{\raggedright\arraybackslash}p{0.09\textwidth}|>{\raggedright\arraybackslash}p{0.54\textwidth}|}
\hline
\textbf{Actor} & \textbf{Use Case} & \textbf{Type} & \textbf{Description} \\ \hline
\endhead
University Student & UC-15 Create Material Post & Primary & Creates and submits a material post with title, description, files, and subject tags. \\ \hline
\end{longtable}
\endgroup
\par\medskip

\hypertarget{sec:a3-2-generalization-and-specialization}{}
\paragraph*{A3.2 Generalization and Specialization}
\par\medskip

\begingroup
\small
\renewcommand{\baselinestretch}{1.0}\selectfont
\renewcommand{\arraystretch}{1.15}
\setlength{\tabcolsep}{6pt}
\begin{longtable}{|>{\raggedright\arraybackslash}p{0.09\textwidth}|>{\raggedright\arraybackslash}p{0.09\textwidth}|>{\raggedright\arraybackslash}p{0.09\textwidth}|>{\raggedright\arraybackslash}p{0.65\textwidth}|}
\hline
\textbf{Child} & \textbf{Parent} & \textbf{Type} & \textbf{Inherited Behavior or Association} \\ \hline
\endhead
None & None & None & No actor generalization or specialization is named for US-15 in the specified proposals. \\ \hline
\end{longtable}
\endgroup
\par\medskip

\hypertarget{sec:a3-3-include}{}
\paragraph*{A3.3 Include}
\par\medskip

\begingroup
\small
\renewcommand{\baselinestretch}{1.0}\selectfont
\renewcommand{\arraystretch}{1.15}
\setlength{\tabcolsep}{6pt}
\begin{longtable}{|>{\raggedright\arraybackslash}p{0.10\textwidth}|>{\raggedright\arraybackslash}p{0.11\textwidth}|>{\raggedright\arraybackslash}p{0.71\textwidth}|}
\hline
\textbf{Base Use Case} & \textbf{Included Use Case} & \textbf{Meaning} \\ \hline
\endhead
None & None & No distinct included use case is named for US-15 in the specified proposals; no new I-xx ID is introduced. \\ \hline
\end{longtable}
\endgroup
\par\medskip

\hypertarget{sec:a3-4-extend}{}
\paragraph*{A3.4 Extend}
\par\medskip

\begingroup
\small
\renewcommand{\baselinestretch}{1.0}\selectfont
\renewcommand{\arraystretch}{1.15}
\setlength{\tabcolsep}{6pt}
\begin{longtable}{|>{\raggedright\arraybackslash}p{0.10\textwidth}|>{\raggedright\arraybackslash}p{0.09\textwidth}|>{\raggedright\arraybackslash}p{0.09\textwidth}|>{\raggedright\arraybackslash}p{0.62\textwidth}|}
\hline
\textbf{Extending Use Case} & \textbf{Base Use Case} & \textbf{Condition} & \textbf{Added Behavior} \\ \hline
\endhead
None & None & None & No distinct extending use case is named for US-15 in the specified proposals; no new E-xx ID is introduced. \\ \hline
\end{longtable}
\endgroup
\par\medskip

\hypertarget{sec:a4-model-notes}{}
\addcontentsline{toc}{subsubsection}{A4. Model Notes}
\subsubsection*{A4. Model Notes}
\par\medskip

\begin{enumerate}
\item UC-15 covers only the Material Post path; Discussion Post and Functional Post are separate paths on the shared Create Post page.
\item The diagram intentionally shows only the source-backed University Student actor and UC-15; no supporting system actor is named for this story.
\item Forum context is university, major, and course as required by \texttt{AC-15.3} and \texttt{FR-15.4}.
\item The optional library listing is an outcome of publication, not a separate use case; priority, dedicated route, accepted file formats, and maximum file sizes are not specified in proposals.
\end{enumerate}
\par\medskip

\hypertarget{sec:part-b-use-case-specification}{}
\addcontentsline{toc}{subsection}{Part B --- Use Case Specification}
\subsection*{Part B --- Use Case Specification}
\par\medskip

\hypertarget{sec:b1-brief-description}{}
\addcontentsline{toc}{subsubsection}{B1. Brief Description}
\subsubsection*{B1. Brief Description}
\par\medskip

A university student creates a material post in a selected forum to share learning resources with other students. The post supports a title, description, files, and subject tags, and uploaded files are validated for type and size before acceptance.
\par\medskip

The material is associated with the selected university, major, and course. Published materials are accessible to users who can view the forum and may also appear in the learning resource library when applicable.
\par\medskip

\hypertarget{sec:b2-preconditions}{}
\addcontentsline{toc}{subsubsection}{B2. Preconditions}
\subsubsection*{B2. Preconditions}
\par\medskip

\begin{enumerate}
\item The student is authenticated with an active session.
\item The student has selected a valid university, major, and course forum context.
\item The student is permitted to post materials in that forum.
\item The student has one or more files intended for the material post.
\end{enumerate}
\par\medskip

\hypertarget{sec:b3-postconditions}{}
\addcontentsline{toc}{subsubsection}{B3. Postconditions}
\subsubsection*{B3. Postconditions}
\par\medskip

\begin{enumerate}
\item \textbf{Success:} A material post is stored with its title, description, validated files, subject tags, author, and forum context.
\item \textbf{Success:} The material post appears in the corresponding forum and is accessible to users who can view the forum.
\item \textbf{Success:} The material may also appear in the learning resource library when applicable.
\item \textbf{Failure:} No material post is published, and the student receives validation, permission, or system feedback.
\end{enumerate}
\par\medskip

\hypertarget{sec:b4-main-success-scenario-basic-flow}{}
\addcontentsline{toc}{subsubsection}{B4. Main Success Scenario (Basic Flow)}
\subsubsection*{B4. Main Success Scenario (Basic Flow)}
\par\medskip

\begin{enumerate}
\item The student selects \textbf{Create Post} from the selected forum.
\item The system shows the active university, major, and course context.
\item The student selects \textbf{Material Post}.
\item The student enters a title and description and adds subject tags.
\item The student selects one or more files through the material file upload control.
\item The system validates file type and size and shows the upload preview.
\item The student reviews the access and library-listing notices.
\item The student submits the post.
\item The frontend validates the material fields and uploads before submission.
\item The backend verifies authentication, forum access, permissions, content, and files.
\item The system stores and publishes the material post and returns the new post result.
\item The post appears in the corresponding forum post list and can be opened in post detail.
\end{enumerate}
\par\medskip

\hypertarget{sec:b5-alternative-flows}{}
\addcontentsline{toc}{subsubsection}{B5. Alternative Flows}
\subsubsection*{B5. Alternative Flows}
\par\medskip

Unless otherwise stated, an alternative flow resumes at the next step of the basic flow.
\par\medskip

\hypertarget{sec:a1-preview-before-submission}{}
\paragraph*{A1. Preview before submission}
\par\medskip

At basic flow step 7, the student previews the title, description, tags, and files. Invalid fields are corrected before the flow resumes at step 8.
\par\medskip

\hypertarget{sec:a2-library-listing-not-applicable}{}
\paragraph*{A2. Library listing not applicable}
\par\medskip

At basic flow step 11, the backend does not include the material in the learning resource library. The forum post is still published because library listing is conditional under \texttt{AC-15.5} and \texttt{FR-15.6}.
\par\medskip

\hypertarget{sec:a3-cancel-creation}{}
\paragraph*{A3. Cancel creation}
\par\medskip

The student leaves the editor before submitting. No material post is published; any retained draft behavior is not specified in proposals.
\par\medskip

\hypertarget{sec:b6-exception-flows}{}
\addcontentsline{toc}{subsubsection}{B6. Exception Flows}
\subsubsection*{B6. Exception Flows}
\par\medskip

\hypertarget{sec:e1-unsupported-or-oversized-file}{}
\paragraph*{E1. Unsupported or oversized file}
\par\medskip

The system rejects the file, identifies the violated type or size constraint, and keeps the editor open. Concrete accepted formats and size limits are not specified in proposals.
\par\medskip

\hypertarget{sec:e2-invalid-post-content-or-post-type}{}
\paragraph*{E2. Invalid post content or post type}
\par\medskip

The system rejects the submission with \texttt{400} and \texttt{INVALID\_POST\_TYPE} (\texttt{48003}) or \texttt{INVALID\_POST\_CONTENT} (\texttt{48004}) and allows the student to correct the material fields.
\par\medskip

\hypertarget{sec:e3-unauthorized-or-forbidden-publication}{}
\paragraph*{E3. Unauthorized or forbidden publication}
\par\medskip

The system rejects the request with \texttt{401}, or with \texttt{403} and \texttt{FORBIDDEN} (\texttt{48010}), when authentication or forum permission is absent.
\par\medskip

\hypertarget{sec:e4-rate-limit-reached}{}
\paragraph*{E4. Rate limit reached}
\par\medskip

The system rejects the submission and asks the student to retry later. \texttt{NFR-SEC-9} requires rate limiting, but no rate-limit status or error code is specified in proposals.
\par\medskip

\hypertarget{sec:e5-server-or-database-failure}{}
\paragraph*{E5. Server or database failure}
\par\medskip

The system returns \texttt{500} with \texttt{DATABASE\_ERROR} or \texttt{UNKNOWN\_ERROR} (\texttt{45000}), publishes no partial post, and offers retry where supported.
\par\medskip

\hypertarget{sec:b7-business-rules}{}
\addcontentsline{toc}{subsubsection}{B7. Business Rules}
\subsubsection*{B7. Business Rules}
\par\medskip

\begin{enumerate}
\item \texttt{AC-15.1} / \texttt{FR-15.2}: The material post supports a title, description, files, and subject tags.
\item \texttt{AC-15.2} / \texttt{FR-15.3}: Uploaded files are validated for type and size before acceptance.
\item \texttt{AC-15.3} / \texttt{FR-15.4}: The material is associated with the selected university, major, and course context.
\item \texttt{AC-15.4} / \texttt{FR-15.5}: Published materials are accessible to users who can view the forum.
\item \texttt{AC-15.5} / \texttt{FR-15.6}: The material may also appear in the learning resource library when applicable.
\item \texttt{FR-15.1}: The system allows a user to create a material post in a selected forum.
\item Authentication, authorization, and server-side permission checks apply before publication (\texttt{NFR-SEC-1}, \texttt{NFR-SEC-6}).
\item Privacy settings apply consistently to forum materials (\texttt{NFR-PRIV-7}), and shared-resource access controls apply where configured (\texttt{NFR-PRIV-5}).
\end{enumerate}
\par\medskip

\hypertarget{sec:b8-special-requirements}{}
\addcontentsline{toc}{subsubsection}{B8. Special Requirements}
\subsubsection*{B8. Special Requirements}
\par\medskip

\begingroup
\small
\renewcommand{\baselinestretch}{1.0}\selectfont
\renewcommand{\arraystretch}{1.15}
\setlength{\tabcolsep}{6pt}
\begin{longtable}{|>{\raggedright\arraybackslash}p{0.13\textwidth}|>{\raggedright\arraybackslash}p{0.80\textwidth}|}
\hline
\textbf{ID} & \textbf{Requirement} \\ \hline
\endhead
NFR-PERF-4 & A new post must be reflected in the relevant forum view within 2 seconds. \\ \hline
NFR-SEC-1 & Authentication is required before accessing protected posting features. \\ \hline
NFR-SEC-5 & Uploaded files must be validated for file type and file size before acceptance. \\ \hline
NFR-SEC-6 & Server-side authorization checks must prevent unauthorized posting. \\ \hline
NFR-SEC-9 & Post creation must be rate-limited to reduce spam and abuse. \\ \hline
NFR-PRIV-5, NFR-PRIV-7 & Users control access to shared resources, and privacy settings apply consistently to forum posts, comments, member lists, and search results. \\ \hline
NFR-SCAL-4 & The architecture supports at least 100,000 uploaded documents and learning resources. \\ \hline
Design constraint & Accepted file formats, maximum file sizes, and upload limits are not specified in proposals. \\ \hline
\end{longtable}
\endgroup
\par\medskip

\hypertarget{sec:b9-extension-points}{}
\addcontentsline{toc}{subsubsection}{B9. Extension Points}
\subsubsection*{B9. Extension Points}
\par\medskip

\begingroup
\small
\renewcommand{\baselinestretch}{1.0}\selectfont
\renewcommand{\arraystretch}{1.15}
\setlength{\tabcolsep}{6pt}
\begin{longtable}{|>{\raggedright\arraybackslash}p{0.38\textwidth}|>{\raggedright\arraybackslash}p{0.55\textwidth}|}
\hline
\textbf{Extension Point} & \textbf{Location in Flow of Events} \\ \hline
\endhead
Select post type & Basic flow step 3 \\ \hline
Add subject tags & Basic flow step 4 \\ \hline
File upload and validation & Basic flow steps 5-6; exception E1 \\ \hline
Preview and library eligibility & Basic flow steps 7 and 11; alternatives A1-A2 \\ \hline
Publication failure recovery & Basic flow step 11; exceptions E2-E5 \\ \hline
\end{longtable}
\endgroup
\par\medskip

\hypertarget{sec:b10-interface-and-data-notes}{}
\addcontentsline{toc}{subsubsection}{B10. Interface and Data Notes}
\subsubsection*{B10. Interface and Data Notes}
\par\medskip

\begin{itemize}
\item \textbf{Screen/flow path:} \texttt{screen\_flow.md} gives \texttt{Forum -$>$ University -$>$ Major -$>$ Course -$>$ Post} and \texttt{Forum -$>$ Create Post -$>$ Post}. \texttt{feature\_flow.md} gives \texttt{Create Post -$>$ Material Post}, with \texttt{Material Post -$>$ Files / Title / Description}, and all post types converging on \textbf{Submit}.
\item \textbf{Editor fields:} \texttt{create\_post.md} recommends \texttt{MaterialEditorForm}, \texttt{PostTitleInput}, \texttt{MaterialDescriptionInput}, \texttt{SubjectTagInput}, \texttt{MaterialFileDropzone}, \texttt{MaterialFilePreviewPanel}, \texttt{FileValidationErrorText}, \texttt{PostPreviewPanel}, \texttt{CreatePostValidationSummary}, and \texttt{SubmitPostButton}.
\item \textbf{Exact recommended UI copy:} \texttt{MaterialFileDropzone} shows “Choose files or drag them here”; \texttt{MaterialAccessNotice} shows “Published materials are accessible to users who can view this forum.”; \texttt{LibraryListingNotice} shows “This material may also appear in the Learning Resource Library.” These strings are recommendations in \texttt{create\_post.md}, not mandated requirement text.
\item \textbf{Forum context and material fields:} \texttt{university\_id: string}, optional \texttt{major\_id?: string}, and optional \texttt{course\_id?: string} associate the material with the forum; the backend body accepts \texttt{post\_type: 'material'}, \texttt{title: string}, \texttt{content: string}, optional \texttt{tags?: string[]}, and optional \texttt{files?: File[]}. \texttt{create\_post.md} recommends \texttt{description}, \texttt{subjectTags}, and \texttt{files} as draft-level fields.
\item \textbf{API endpoint:} \texttt{POST /forum/posts} returns \texttt{201} with \texttt{data: \{ post\_id: string, created\_at: number \}}; failures are \texttt{400} with \texttt{(INVALID\_POST\_TYPE | INVALID\_POST\_CONTENT)[]}, \texttt{401}, \texttt{403} with \texttt{(FORBIDDEN)[]}, and \texttt{500} with \texttt{(DATABASE\_ERROR | UNKNOWN\_ERROR)[]}. \texttt{studitopia.create\_forum\_post(pauthor\_id: uuid, puniversity\_id: uuid, pmajor\_id: uuid, pcourse\_id: uuid, ppost\_type: string, ptitle: string, pcontent: string, ptags: string[], pfiles: jsonb)} returns \texttt{\{post\_id: uuid, created\_at: timestamp\}} and declares \texttt{UNKNOWN\_ERROR: 45000}, \texttt{INVALID\_POST\_TYPE: 48003}, \texttt{INVALID\_POST\_CONTENT: 48004}, and \texttt{FORBIDDEN: 48010}.
\item \textbf{Post-publication surface and gaps:} \texttt{list\_post.md} requires the summary to show title, author, post type, timestamp, vote score, comment count, solved status, and tags, with attachments in post detail; dedicated route, accepted file formats, maximum file sizes, and a rate-limit error code are not specified in proposals.
\end{itemize}
\par\medskip

\hypertarget{sec:b11-source-alignment-notes}{}
\addcontentsline{toc}{subsubsection}{B11. Source Alignment Notes}
\subsubsection*{B11. Source Alignment Notes}
\par\medskip

\begin{itemize}
\item UC-04 covers only the Discussion Post path; UC-15 shares the Create Post entry point and \texttt{POST /forum/posts} contract but does not inherit LaTeX, images, or question/general category. UC-05 covers comment and answer creation after publication and does not contribute steps to material creation.
\item \texttt{userstory.md} states “Create Material Post,” while \texttt{functional\_requirement.md} states “Create a material post, so that I can share learning resources with other students.” This document uses the exact \texttt{userstory.md} title; \texttt{API.md} has one \texttt{post\_type} enum for \texttt{discussion | material | functional}, and UC-15 constrains its flow to \texttt{material}.
\item \texttt{create\_post.md} is explicitly a UI/UX recommendation; its component names and quoted messages are not mandated by FR/AC text. The optional library listing in \texttt{AC-15.5} / \texttt{FR-15.6} is conditional, and no dedicated library-publication endpoint is specified in the read sources.
\item No distinct included, extending, or specialized use case is named for US-15, so no new I-xx/E-xx/G-xx ID is introduced.
\end{itemize}
\par\medskip

\par\medskip
\hypertarget{sec:ui-prototype-specimens}{}
\addcontentsline{toc}{subsection}{UI Prototype Specimens}
\subsection*{UI Prototype Specimens}
\par\medskip
\begin{figure}[H]
\centering
\includegraphics[width=0.75\textwidth,height=0.42\textheight,keepaspectratio]{"../../prototype/create_post_material.png"}
\par\vspace{2pt}
{\footnotesize \textbf{UI Prototype:} Create post - Material}
\end{figure}
\par\smallskip

\pagebreak
\hypertarget{sec:uc-16-create-functional-post}{}
\addcontentsline{toc}{section}{UC-16 --- Create Functional Post}
\section*{UC-16 --- Create Functional Post}
\par\medskip

\textbf{Version:} \texttt{1.0}
\par\medskip

\textbf{Performed by:} [Nguyen Thanh An] | \textbf{Edited by:} [Nguyen Thanh An]
\par\medskip

\begingroup
\small
\renewcommand{\baselinestretch}{1.0}\selectfont
\renewcommand{\arraystretch}{1.15}
\setlength{\tabcolsep}{6pt}
\begin{longtable}{|>{\raggedright\arraybackslash}p{0.09\textwidth}|>{\raggedright\arraybackslash}p{0.84\textwidth}|}
\hline
\textbf{Field} & \textbf{Value} \\ \hline
\endhead
Use Case ID & UC-16 \\ \hline
Name & Create Functional Post \\ \hline
Primary Actor & University Student \\ \hline
Supporting Actors & None \\ \hline
Feature Group & Post Management \\ \hline
Priority & Not explicitly stated \\ \hline
Related User Stories & US-16 \\ \hline
Related Requirements & FR-16.1 to FR-16.6; NFR-PERF-4, NFR-SEC-1, NFR-SEC-6, NFR-SEC-8, NFR-SEC-9, NFR-PRIV-7 \\ \hline
Source Documents & \texttt{proposals/userstory.md} (section 5.3); \texttt{proposals/functional\_requirement.md} (US-16); \texttt{proposals/create\_post.md}; \texttt{proposals/list\_post.md}; \texttt{proposals/post\_management.md}; \texttt{proposals/API.md}; \texttt{proposals/response\_code.md}; \texttt{proposals/non\_functional\_requirement.md}; \texttt{proposals/feature\_flow.md}; \texttt{proposals/screen\_flow.md} \\ \hline
\end{longtable}
\endgroup
\par\medskip

\hypertarget{sec:part-a-use-case-model}{}
\addcontentsline{toc}{subsection}{Part A --- Use Case Model}
\subsection*{Part A --- Use Case Model}
\par\medskip

\hypertarget{sec:a1-focused-use-case-diagram}{}
\addcontentsline{toc}{subsubsection}{A1. Focused Use Case Diagram}
\subsubsection*{A1. Focused Use Case Diagram}
\par\medskip

\begin{center}
\includegraphics[width=0.85\textwidth,height=0.75\textheight,keepaspectratio]{"figures/UC-16-create-functional-post_1.png"}
\end{center}
\par\medskip

\textbf{Diagram notation.} Person shape = human actor; cloud = supporting system; blue = selected use case; solid line = association; labelled solid arrow = generalization; dashed arrow = include or extend.
\par\medskip

\hypertarget{sec:a2-actor-register}{}
\addcontentsline{toc}{subsubsection}{A2. Actor Register}
\subsubsection*{A2. Actor Register}
\par\medskip

\begingroup
\small
\renewcommand{\baselinestretch}{1.0}\selectfont
\renewcommand{\arraystretch}{1.15}
\setlength{\tabcolsep}{6pt}
\begin{longtable}{|>{\raggedright\arraybackslash}p{0.10\textwidth}|>{\raggedright\arraybackslash}p{0.10\textwidth}|>{\raggedright\arraybackslash}p{0.72\textwidth}|}
\hline
\textbf{Actor} & \textbf{Type} & \textbf{Description} \\ \hline
\endhead
University Student & Primary human actor & A university student who creates an interactive academic tool in a selected forum so that other students can practice or understand course content. \\ \hline
\end{longtable}
\endgroup
\par\medskip

\hypertarget{sec:a3-relationships}{}
\addcontentsline{toc}{subsubsection}{A3. Relationships}
\subsubsection*{A3. Relationships}
\par\medskip

\hypertarget{sec:a3-1-associations}{}
\paragraph*{A3.1 Associations}
\par\medskip

\begingroup
\small
\renewcommand{\baselinestretch}{1.0}\selectfont
\renewcommand{\arraystretch}{1.15}
\setlength{\tabcolsep}{6pt}
\begin{longtable}{|>{\raggedright\arraybackslash}p{0.09\textwidth}|>{\raggedright\arraybackslash}p{0.15\textwidth}|>{\raggedright\arraybackslash}p{0.09\textwidth}|>{\raggedright\arraybackslash}p{0.57\textwidth}|}
\hline
\textbf{Actor} & \textbf{Use Case} & \textbf{Type} & \textbf{Description} \\ \hline
\endhead
University Student & UC-16 Create Functional Post & Primary & Configures and publishes a functional post containing a quiz, prerequisite tree, or dynamic computation. \\ \hline
\end{longtable}
\endgroup
\par\medskip

\hypertarget{sec:a3-2-generalization-and-specialization}{}
\paragraph*{A3.2 Generalization and Specialization}
\par\medskip

\begingroup
\small
\renewcommand{\baselinestretch}{1.0}\selectfont
\renewcommand{\arraystretch}{1.15}
\setlength{\tabcolsep}{6pt}
\begin{longtable}{|>{\raggedright\arraybackslash}p{0.09\textwidth}|>{\raggedright\arraybackslash}p{0.09\textwidth}|>{\raggedright\arraybackslash}p{0.09\textwidth}|>{\raggedright\arraybackslash}p{0.65\textwidth}|}
\hline
\textbf{Child} & \textbf{Parent} & \textbf{Type} & \textbf{Inherited Behavior or Association} \\ \hline
\endhead
None & None & None & No actor generalization or specialization is named for US-16 in the specified proposals. \\ \hline
\end{longtable}
\endgroup
\par\medskip

\hypertarget{sec:a3-3-include}{}
\paragraph*{A3.3 Include}
\par\medskip

\begingroup
\small
\renewcommand{\baselinestretch}{1.0}\selectfont
\renewcommand{\arraystretch}{1.15}
\setlength{\tabcolsep}{6pt}
\begin{longtable}{|>{\raggedright\arraybackslash}p{0.10\textwidth}|>{\raggedright\arraybackslash}p{0.11\textwidth}|>{\raggedright\arraybackslash}p{0.71\textwidth}|}
\hline
\textbf{Base Use Case} & \textbf{Included Use Case} & \textbf{Meaning} \\ \hline
\endhead
None & None & No distinct included use case is named for US-16 in the specified proposals; no new I-xx ID is introduced. \\ \hline
\end{longtable}
\endgroup
\par\medskip

\hypertarget{sec:a3-4-extend}{}
\paragraph*{A3.4 Extend}
\par\medskip

\begingroup
\small
\renewcommand{\baselinestretch}{1.0}\selectfont
\renewcommand{\arraystretch}{1.15}
\setlength{\tabcolsep}{6pt}
\begin{longtable}{|>{\raggedright\arraybackslash}p{0.10\textwidth}|>{\raggedright\arraybackslash}p{0.09\textwidth}|>{\raggedright\arraybackslash}p{0.09\textwidth}|>{\raggedright\arraybackslash}p{0.62\textwidth}|}
\hline
\textbf{Extending Use Case} & \textbf{Base Use Case} & \textbf{Condition} & \textbf{Added Behavior} \\ \hline
\endhead
None & None & None & No distinct extending use case is named for US-16 in the specified proposals; no new E-xx ID is introduced. \\ \hline
\end{longtable}
\endgroup
\par\medskip

\hypertarget{sec:a4-model-notes}{}
\addcontentsline{toc}{subsubsection}{A4. Model Notes}
\subsubsection*{A4. Model Notes}
\par\medskip

\begin{enumerate}
\item UC-16 covers only the Functional Post path; Discussion Post and Material Post are separate paths on the shared Create Post page.
\item The three tool variants are source-backed configurations in \texttt{feature\_flow.md}, \texttt{create\_post.md}, \texttt{FR-16.3} to \texttt{FR-16.5}, and \texttt{AC-16.1} to \texttt{AC-16.3}; they are not separate use cases.
\item \texttt{create\_post.md} recommends drag-and-drop, code, and generate-from-file modes for prerequisite trees; no new included/extending use case is created for those modes.
\item The selected forum context is university, major, and course under \texttt{AC-16.5}; priority, exact code syntax, source-file formats, tree schema details, and numeric limits are not specified in proposals.
\end{enumerate}
\par\medskip

\hypertarget{sec:part-b-use-case-specification}{}
\addcontentsline{toc}{subsection}{Part B --- Use Case Specification}
\subsection*{Part B --- Use Case Specification}
\par\medskip

\hypertarget{sec:b1-brief-description}{}
\addcontentsline{toc}{subsubsection}{B1. Brief Description}
\subsubsection*{B1. Brief Description}
\par\medskip

A university student creates a functional post in a selected forum to help others practice or understand course content. The post contains a title and one interactive academic tool: an immediate quiz, a prerequisite tree, or dynamic computation. The author configures and previews the tool before submission; after publication, other users can interact with the published functional tool from post detail.
\par\medskip

\hypertarget{sec:b2-preconditions}{}
\addcontentsline{toc}{subsubsection}{B2. Preconditions}
\subsubsection*{B2. Preconditions}
\par\medskip

\begin{enumerate}
\item The student is authenticated with an active session.
\item The student has selected a valid forum context.
\item The student is permitted to create posts in that forum.
\item The author has chosen the interactive tool to configure.
\end{enumerate}
\par\medskip

\hypertarget{sec:b3-postconditions}{}
\addcontentsline{toc}{subsubsection}{B3. Postconditions}
\subsubsection*{B3. Postconditions}
\par\medskip

\begin{enumerate}
\item \textbf{Success:} A functional post is stored with its title, selected tool configuration, author, and forum context.
\item \textbf{Success:} Other users can interact with the published tool from post detail, and the post appears in the corresponding forum post list.
\item \textbf{Failure:} No functional post is published, and the author receives validation, permission, or system feedback.
\end{enumerate}
\par\medskip

\hypertarget{sec:b4-main-success-scenario-basic-flow}{}
\addcontentsline{toc}{subsubsection}{B4. Main Success Scenario (Basic Flow)}
\subsubsection*{B4. Main Success Scenario (Basic Flow)}
\par\medskip

\begin{enumerate}
\item The student selects \textbf{Create Post} from the selected forum.
\item The system shows the active university, major, and course context.
\item The student selects \textbf{Functional Post}.
\item The student enters a title and selects one tool: \textbf{Immediate Quiz}, \textbf{Prerequisite Tree}, or \textbf{Dynamic Computation}.
\item The student configures the selected tool and previews it.
\item The system validates the tool configuration and marks the draft \textbf{Ready} when it can be submitted.
\item The student submits the post.
\item The frontend validates required metadata and the active tool configuration.
\item The backend verifies authentication, forum access, permissions, content, and rate limits.
\item The system stores and publishes the functional post and returns the new post result.
\item Other users can open the post and interact with the published tool.
\end{enumerate}
\par\medskip

\hypertarget{sec:b5-alternative-flows}{}
\addcontentsline{toc}{subsubsection}{B5. Alternative Flows}
\subsubsection*{B5. Alternative Flows}
\par\medskip

Unless otherwise stated, an alternative flow resumes at the next step of the basic flow.
\par\medskip

\hypertarget{sec:a1-immediate-quiz}{}
\paragraph*{A1. Immediate quiz}
\par\medskip

At basic flow step 4, the student selects \textbf{Immediate Quiz}. The quiz supports multiple-choice or single-choice questions, options, correct answers, and scores under \texttt{AC-16.2} and \texttt{FR-16.3}.
\par\medskip

\hypertarget{sec:a2-prerequisite-tree}{}
\paragraph*{A2. Prerequisite tree}
\par\medskip

At basic flow step 4, the student selects \textbf{Prerequisite Tree}. The tree can be created manually or generated from a source file under \texttt{AC-16.3} and \texttt{FR-16.4}; \texttt{create\_post.md} recommends \textbf{Drag and drop}, \textbf{Code}, and \textbf{Generate from file} modes.
\par\medskip

\hypertarget{sec:a3-dynamic-computation}{}
\paragraph*{A3. Dynamic computation}
\par\medskip

At basic flow step 4, the student selects \textbf{Dynamic Computation}. The author configures inputs, formula/configuration, and outputs under \texttt{FR-16.5}; no computation schema is specified in proposals.
\par\medskip

\hypertarget{sec:a4-cancel-creation}{}
\paragraph*{A4. Cancel creation}
\par\medskip

The student leaves the editor before submitting. No functional post is published; retained draft behavior is not specified in proposals.
\par\medskip

\hypertarget{sec:b6-exception-flows}{}
\addcontentsline{toc}{subsubsection}{B6. Exception Flows}
\subsubsection*{B6. Exception Flows}
\par\medskip

\hypertarget{sec:e1-invalid-or-incomplete-tool-configuration}{}
\paragraph*{E1. Invalid or incomplete tool configuration}
\par\medskip

The system shows validation feedback and keeps the editor open. A quiz with invalid answers, a tree with no nodes, or an invalid computation configuration cannot be published.
\par\medskip

\hypertarget{sec:e2-prerequisite-tree-needs-review}{}
\paragraph*{E2. Prerequisite tree needs review}
\par\medskip

A generated tree with uncertain nodes or edges is marked \textbf{Needs review} (\texttt{NeedsReview}) and must be confirmed by the author before publication. A draft with missing or invalid nodes/edges is \textbf{Incomplete} and blocked.
\par\medskip

\hypertarget{sec:e3-invalid-post-content-or-post-type}{}
\paragraph*{E3. Invalid post content or post type}
\par\medskip

The system rejects the submission with \texttt{400} and \texttt{INVALID\_POST\_TYPE} (\texttt{48003}) or \texttt{INVALID\_POST\_CONTENT} (\texttt{48004}). No tool-specific error code is defined in \texttt{API.md}.
\par\medskip

\hypertarget{sec:e4-unauthorized-or-forbidden-publication}{}
\paragraph*{E4. Unauthorized or forbidden publication}
\par\medskip

The system rejects the request with \texttt{401}, or with \texttt{403} and \texttt{FORBIDDEN} (\texttt{48010}), when authentication or forum permission is absent.
\par\medskip

\hypertarget{sec:e5-rate-limit-reached}{}
\paragraph*{E5. Rate limit reached}
\par\medskip

The system rejects the submission and asks the student to retry later. \texttt{NFR-SEC-9} requires rate limiting, but no rate-limit status or error code is specified in proposals.
\par\medskip

\hypertarget{sec:e6-server-or-database-failure}{}
\paragraph*{E6. Server or database failure}
\par\medskip

The system returns \texttt{500} with \texttt{DATABASE\_ERROR} or \texttt{UNKNOWN\_ERROR} (\texttt{45000}), publishes no partial tool, and offers retry where supported.
\par\medskip

\hypertarget{sec:b7-business-rules}{}
\addcontentsline{toc}{subsubsection}{B7. Business Rules}
\subsubsection*{B7. Business Rules}
\par\medskip

\begin{enumerate}
\item \texttt{AC-16.1} / \texttt{FR-16.2}: The post supports a title and an interactive tool.
\item \texttt{AC-16.1} / \texttt{FR-16.3} to \texttt{FR-16.5}: The tool may be an immediate quiz, prerequisite tree, or dynamic computation.
\item \texttt{AC-16.2} / \texttt{FR-16.3}: Quizzes support multiple-choice or single-choice questions, options, answers, and scores.
\item \texttt{AC-16.3} / \texttt{FR-16.4}: Prerequisite trees can be created manually or generated from a source file.
\item \texttt{AC-16.4} / \texttt{FR-16.6}: Other users can interact with the published functional tool.
\item \texttt{AC-16.5}: The functional post is associated with the selected forum context. A prerequisite-tree node represents a course, topic, concept, skill, or milestone; a directed edge runs from a prerequisite node to a dependent node.
\item Tree edges must reference existing nodes, a node cannot be its own prerequisite, and cycles must be flagged before publish; these edge rules are recommendations recorded in \texttt{create\_post.md}. Drag-and-drop and code authoring must produce the same recommended \texttt{treeConfig}, and code must be applied to preview before submission.
\item Generated trees remain editable drafts until the author reviews and confirms them; readiness states are \texttt{Incomplete}, \texttt{NeedsReview} (shown as \textbf{Needs review}), and \texttt{Ready}. No tree-specific API error code is defined; invalid tool content is mapped through generic validation feedback such as \texttt{INVALID\_POST\_CONTENT}.
\end{enumerate}
\par\medskip

\hypertarget{sec:b8-special-requirements}{}
\addcontentsline{toc}{subsubsection}{B8. Special Requirements}
\subsubsection*{B8. Special Requirements}
\par\medskip

\begingroup
\small
\renewcommand{\baselinestretch}{1.0}\selectfont
\renewcommand{\arraystretch}{1.15}
\setlength{\tabcolsep}{6pt}
\begin{longtable}{|>{\raggedright\arraybackslash}p{0.10\textwidth}|>{\raggedright\arraybackslash}p{0.83\textwidth}|}
\hline
\textbf{ID} & \textbf{Requirement} \\ \hline
\endhead
NFR-PERF-4 & A new post must be reflected in the relevant forum view within 2 seconds. \\ \hline
NFR-SEC-1 & Authentication is required before accessing protected posting features. \\ \hline
NFR-SEC-6 & Server-side authorization checks must prevent unauthorized posting. \\ \hline
NFR-SEC-8 & The system must protect against SQL injection, XSS, and CSRF. \\ \hline
NFR-SEC-9 & Post creation must be rate-limited to reduce spam and abuse. \\ \hline
NFR-PRIV-7 & Privacy settings apply consistently to forum posts, comments, member lists, and search results. \\ \hline
Design constraint & Exact tree schema, code syntax, source-file formats, generation rules, and numeric limits are not specified in proposals; the published tool viewer should use the same reviewed configuration as the author preview. \\ \hline
\end{longtable}
\endgroup
\par\medskip

\hypertarget{sec:b9-extension-points}{}
\addcontentsline{toc}{subsubsection}{B9. Extension Points}
\subsubsection*{B9. Extension Points}
\par\medskip

\begingroup
\small
\renewcommand{\baselinestretch}{1.0}\selectfont
\renewcommand{\arraystretch}{1.15}
\setlength{\tabcolsep}{6pt}
\begin{longtable}{|>{\raggedright\arraybackslash}p{0.44\textwidth}|>{\raggedright\arraybackslash}p{0.49\textwidth}|}
\hline
\textbf{Extension Point} & \textbf{Location in Flow of Events} \\ \hline
\endhead
Select interactive tool & Basic flow step 4; alternatives A1-A3 \\ \hline
Quiz and tree configuration & Basic flow step 5; alternatives A1-A2 \\ \hline
Generated tree review & Basic flow step 5; exception E2 \\ \hline
Dynamic computation configuration & Basic flow step 5; alternative A3 \\ \hline
Publication failure recovery & Basic flow step 10; exceptions E3-E6 \\ \hline
\end{longtable}
\endgroup
\par\medskip

\hypertarget{sec:b10-interface-and-data-notes}{}
\addcontentsline{toc}{subsubsection}{B10. Interface and Data Notes}
\subsubsection*{B10. Interface and Data Notes}
\par\medskip

\begin{itemize}
\item \textbf{Screen/flow path:} \texttt{screen\_flow.md} gives \texttt{Forum -$>$ University -$>$ Major -$>$ Course -$>$ Post} and \texttt{Forum -$>$ Create Post -$>$ Post}. \texttt{feature\_flow.md} gives \texttt{Create Post -$>$ Functional Post -$>$ Title / Quiz / Prerequisite Tree / Dynamic Computation}, \texttt{Prerequisite Tree -$>$ Auto apply / Simulate by click}, and all post types converging on \textbf{Submit}.
\item \textbf{Editor components and quiz:} \texttt{create\_post.md} recommends \texttt{FunctionalEditorForm}, \texttt{InteractiveToolSelector}, \texttt{QuizBuilder}, \texttt{PrerequisiteTreeBuilder}, \texttt{DynamicComputationConfigurator}, \texttt{TreePreviewPanel}, \texttt{TreeValidationSummary}, and \texttt{SubmitPostButton}. \texttt{QuizBuilder} covers questions, multiple-choice or single-choice type, options, correct-answer markers, and scores; \texttt{QuizPreviewPanel} tests answer submission.
\item \textbf{Prerequisite-tree authoring and rules:} \texttt{TreeModeSelector} offers \textbf{Drag and drop}, \textbf{Code}, and \textbf{Generate from file}. Drag-and-drop uses \texttt{AddPrerequisiteNodeButton}, \texttt{ManualPrerequisiteNodeEditor}, and \texttt{PrerequisiteEdgeBuilder}; code uses \texttt{CodeTreeEditor}, \texttt{ApplyCodeToPreviewButton}, \texttt{CodeValidationPanel}, and optional \texttt{SyncCodeFromCanvasButton}; generated files use \texttt{TreeSourceFileDropzone}, \texttt{GenerateTreeFromFileButton}, and \texttt{GeneratedTreeDraftReview}. A node is a course, topic, concept, skill, or milestone; each directed edge links a prerequisite to a dependent node. Nodes with no incoming prerequisite edges are starting nodes; edges must reference existing nodes; self-dependencies are blocked; cycles are flagged for removal or reordering.
\item \textbf{Simulation, states, and data:} \texttt{AutoApplyToggle} and \texttt{SimulateByClickButton} implement flow-derived \textbf{Auto apply} and \textbf{Simulate by click} behavior. \texttt{TreePublishReadinessBadge} shows \textbf{Ready}, \textbf{Incomplete}, or \textbf{Needs review} (\texttt{treeConfig.validationState} uses \texttt{NeedsReview}), with no numeric bounds specified. \texttt{toolType}, \texttt{quizConfig}, \texttt{treeConfig}, \texttt{computationConfig}, \texttt{treeConfig.nodes[]}, \texttt{treeConfig.edges[]}, \texttt{edges[].fromNodeId}, \texttt{edges[].toNodeId}, \texttt{treeConfig.sourceCode}, \texttt{treeConfig.generationSource}, and \texttt{treeConfig.simulation} are recommendations; \texttt{API.md} has no tree-specific payload field.
\item \textbf{API endpoint and failures:} \texttt{POST /forum/posts} accepts \texttt{university\_id: string}, optional \texttt{major\_id?: string}, optional \texttt{course\_id?: string}, \texttt{post\_type: 'functional'}, \texttt{title: string}, \texttt{content: string}, optional \texttt{tags?: string[]}, and optional \texttt{files?: File[]}; success is \texttt{201} with \texttt{data: \{ post\_id: string, created\_at: number \}}. Failures are \texttt{400} with \texttt{(INVALID\_POST\_TYPE | INVALID\_POST\_CONTENT)[]}, \texttt{401}, \texttt{403} with \texttt{(FORBIDDEN)[]}, and \texttt{500} with \texttt{(DATABASE\_ERROR | UNKNOWN\_ERROR)[]}; database codes are \texttt{UNKNOWN\_ERROR: 45000}, \texttt{INVALID\_POST\_TYPE: 48003}, \texttt{INVALID\_POST\_CONTENT: 48004}, and \texttt{FORBIDDEN: 48010}.
\item \textbf{Reader tools and gaps:} \texttt{list\_post.md} recommends \texttt{QuizRunner}, \texttt{PrerequisiteTreeViewer}, and \texttt{DynamicComputationRunner} in the published functional tool panel; route, accepted tree source-file formats, code syntax, detailed node/edge schema, generation algorithm, and numeric limits are not specified in proposals.
\end{itemize}
\par\medskip

\hypertarget{sec:b11-source-alignment-notes}{}
\addcontentsline{toc}{subsubsection}{B11. Source Alignment Notes}
\subsubsection*{B11. Source Alignment Notes}
\par\medskip

\begin{itemize}
\item UC-04 covers the Discussion Post path on the same Create Post page; UC-16 does not import discussion fields, LaTeX, images, or question/general category. UC-05 covers comments and answers on an existing post and does not contribute steps to functional-tool creation or publication.
\item \texttt{userstory.md} names the story “Create Functional Post” and says “interactive academic tool,” while \texttt{functional\_requirement.md} says “practice course content”; this document uses the exact \texttt{userstory.md} title. \texttt{feature\_flow.md} names only Quiz and Prerequisite Tree under the creation branch; \texttt{create\_post.md} and \texttt{FR-16.5} preserve Dynamic Computation as an explicit tool.
\item The drag-and-drop and code authoring modes are user-direction recommendations recorded in \texttt{create\_post.md}; manual creation and generate-from-file are the explicit \texttt{AC-16.3} / \texttt{FR-16.4} requirement. \texttt{API.md} accepts generic \texttt{content} and \texttt{files} fields and has no tree-specific payload or error contract; readiness states and \texttt{treeConfig} remain UI recommendations.
\item No distinct included, extending, or specialized use case is named for US-16 in the specified proposals, so no new I-xx/E-xx/G-xx ID is introduced.
\end{itemize}
\par\medskip

\par\medskip
\hypertarget{sec:ui-prototype-specimens}{}
\addcontentsline{toc}{subsection}{UI Prototype Specimens}
\subsection*{UI Prototype Specimens}
\par\medskip
\begin{figure}[H]
\centering
\includegraphics[width=0.75\textwidth,height=0.42\textheight,keepaspectratio]{"../../prototype/create_post_functional_quiz.png"}
\par\vspace{2pt}
{\footnotesize \textbf{UI Prototype:} Create post - Functional - Quiz}
\end{figure}
\par\smallskip
\begin{figure}[H]
\centering
\includegraphics[width=0.75\textwidth,height=0.42\textheight,keepaspectratio]{"../../prototype/create_post_prerequisite_tree_code.png"}
\par\vspace{2pt}
{\footnotesize \textbf{UI Prototype:} Create post - Prerequisite tree - Code}
\end{figure}
\par\smallskip
\begin{figure}[H]
\centering
\includegraphics[width=0.75\textwidth,height=0.42\textheight,keepaspectratio]{"../../prototype/create_post_prerequisite_tree_drag_and_drop.png"}
\par\vspace{2pt}
{\footnotesize \textbf{UI Prototype:} Create post - Prerequisite tree - Drag and drop}
\end{figure}
\par\smallskip
\begin{figure}[H]
\centering
\includegraphics[width=0.75\textwidth,height=0.42\textheight,keepaspectratio]{"../../prototype/create_post_prerequisite_tree_generate_from_file.png"}
\par\vspace{2pt}
{\footnotesize \textbf{UI Prototype:} Create post - Prerequisite tree - Generate from file}
\end{figure}
\par\smallskip
\begin{figure}[H]
\centering
\includegraphics[width=0.75\textwidth,height=0.42\textheight,keepaspectratio]{"../../prototype/prerequisite_tree_validation_and_recovery_states.png"}
\par\vspace{2pt}
{\footnotesize \textbf{UI Prototype:} Prerequisite tree - Validation and recovery states}
\end{figure}
\par\smallskip
\begin{figure}[H]
\centering
\includegraphics[width=0.75\textwidth,height=0.42\textheight,keepaspectratio]{"../../prototype/published_prerequisite_tree_reader_simulation.png"}
\par\vspace{2pt}
{\footnotesize \textbf{UI Prototype:} Published prerequisite tree - Reader simulation}
\end{figure}
\par\smallskip

\end{document}
